\documentclass[11pt,letterpaper]{article}
\usepackage[margin=1.02in,top=.90in,bottom=.94in]{geometry}
\usepackage{amsmath,amssymb,amsthm,booktabs,xcolor}
\usepackage[protrusion=true,expansion=false]{microtype}
\definecolor{ink}{rgb}{.08,.22,.35}
\usepackage[colorlinks=true,linkcolor=ink,citecolor=ink,urlcolor=ink,pdftitle={Distinguishing Dirichlet characters from low zeros}]{hyperref}
\usepackage{fancyhdr}
\pagestyle{fancy}\fancyhf{}\fancyhead[L]{\small\itshape Distinguishing Dirichlet characters from low zeros}\fancyhead[R]{\small Research note}\fancyfoot[C]{\thepage}\setlength{\headheight}{14pt}
\newtheorem{theorem}{Theorem}
\newtheorem{lemma}[theorem]{Lemma}
\newtheorem{corollary}[theorem]{Corollary}
\newtheorem{proposition}[theorem]{Proposition}
\newcommand{\C}{\mathbb C}\newcommand{\R}{\mathbb R}\newcommand{\Z}{\mathbb Z}
\newcommand{\F}{\mathcal F}\newcommand{\Pp}{\mathcal P}
\newcommand{\dd}{\,\mathrm d}
\DeclareMathOperator{\sinc}{sinc}
\DeclareMathOperator{\supp}{supp}
\setlength{\emergencystretch}{1.9em}
\raggedbottom
\begin{document}
\thispagestyle{plain}
\begin{center}
{\Large\bfseries Distinguishing Dirichlet characters from low zeros}\par
\vspace{5pt}
{\large A conductor-uniform bound from averaged prime-pair estimates}\par
\vspace{9pt}
{\small 26 November 2026}
\end{center}
\begin{abstract}
Let $Q\ge3$. Using Dirichlet GRH, we prove that the positive zero ordinates up to $C\log^2(2Q)/\log\log(3Q)$, counted with multiplicity, identify every primitive character of conductor at most $Q$, of arbitrary order. At this height distinct characters have $\gg\log(2Q)$ unmatched zero occurrences, even if matching allows a displacement of $c/\log^2(2Q)$. Folding the full signed spectrum by $\gamma\mapsto|\gamma|$ identifies the character up to conjugation. The proof constructs $O(\log Q)$ common prime evaluations which separate all pairs with a fixed margin. A smooth explicit-formula probe reads each evaluation to fixed accuracy. An upper-bound sieve controls the total contribution of nearby prime powers averaged over the probe primes. This permits frequency resolution on the scale of the mean prime spacing. The result answers the finite-height questions for quadratic and complex characters in the ProveIt article.
\end{abstract}

\section{The height and what the data determine}
Let $\mathcal C(Q)$ comprise primitive Dirichlet characters of conductor $q\le Q$, including the character of conductor one for $\zeta$. Characters are identified by their primitive representatives. Put
\[
 L=\log(2Q),\qquad H_Q=\frac{L^2}{\log\log(3Q)} .
\]
All logarithms are natural. Write $m_\chi(\gamma)$ for the multiplicity of the nontrivial zero $1/2+i\gamma$ of $L(s,\chi)$, with value zero at a regular point. All results about zero lists below use the critical-line theorem for Dirichlet $L$-functions. All constants are absolute unless a fixed test function is indicated.

\begin{theorem}[Finite identification]\label{th:main}
There are absolute $C,c>0$ with the following properties.
\begin{enumerate}
\item The multiset of ordinates $0<\gamma\le C H_Q$ determines $\chi\in\mathcal C(Q)$ uniquely.
\item For any two distinct members, even after matching occurrences in these lists at distance at most $c/L^2$, at least $cL$ occurrences in the two lists together remain unmatched.
\item If the data instead consist of the multiset of $|\gamma|$ for all signed ordinates with $|\gamma|\le C H_Q$, they determine the unordered pair $\{\chi,\bar\chi\}$. Conjugate characters give exactly the same folded data at every height.
\end{enumerate}
\end{theorem}

A positive list in (i) keeps only zeros above the real axis. A folded list in (iii) includes contributions from both half-planes, superimposed at the same nonnegative ordinate. Their information content differs for a complex character. Allowing a shared ordinate with different multiplicities to contribute to the discrepancy is essential to the statement of (ii). Primitivity matters because removing finitely many Euler factors leaves the open-strip zeros unchanged. Corollary~\ref{cor:packet} also recovers a bounded number of primitive factors from the zeros of their product, even when the individual zeros have no factor labels.

\clearpage
The following estimate gives the dependence on the height of the observing interval and the total displacement of matched ordinates.

\begin{theorem}[Translated intervals]\label{th:shift}
Let $X=A L^2$, with $A$ a sufficiently large absolute constant. For $U\in\R$ and
\begin{equation}\label{eq:height}
 R\ \ge\ B\left(\frac X{\log X}+\log\log\big(e^e Q(3+|U|)\big)\right),
\end{equation}
where $B$ is an absolute constant, let $I=[U-R,U+R]$. Suppose occurrences from the zero lists in $I$ of two distinct characters in $\mathcal C(Q)$ are partially matched, each occurrence being used at most once. If $E$ is the total number left over and $D$ the sum of the matched distances, then
\begin{equation}\label{eq:transport}
 E+(\log X)D\ \gg\ \frac{R\log X}{\sqrt X}.
\end{equation}
In particular $\sum_{\gamma\in I}|m_{\chi_1}(\gamma)-m_{\chi_2}(\gamma)|\gg R\log X/\sqrt X$, the sum being over distinct ordinates in the union.
\end{theorem}

\paragraph{The two repository inputs.}
The Lean source at commit \texttt{72fead102f66} specifies the Riemann hypothesis for $\zeta$ and the open-strip zero statement for Dirichlet characters, including the primitive case. The proofs are exported in \texttt{grh/src/Main.lean} \cite{lean}. The upstream argument passes from the comparison of corrected lattice covolumes to a one-sided bound for a smoothed difference of prime sums, then deduces nonvanishing to the right of $1/2$. The present proof uses the resulting critical-line location for the original characters and the primitive characters inducing their products. The pole of $\zeta$ is retained explicitly. The upstream repository reports its Lean and independent checker runs. We inspected the source declarations and did not rerun those builds. The new arguments in this note are mathematical proofs outside Lean.

At ProveIt commit \texttt{10802a6b360d}, the section on reconstructing characters contains a compactly supported test isolating each individual prime power and a finite rounding argument \cite{article}. Its Corollary \texttt{rec:cor:quadsep} gives $O(L^2\log\log(3Q))$ for quadratic characters. Its question on the height needed for recovery asks about deleting the factor $\log\log(3Q)$ and about separating complex characters, including conjugates. Theorem~\ref{th:main} improves the stated upper bound by two factors of order $\log\log Q$ and includes unrestricted order. The article is a synthesis of 103 reports, and its README describes their verification scope. No theorem from those reports is needed as a black box here.

\paragraph{Relation to earlier results.}
Banks \cite{banks} obtains unconditionally a discrepancy for distinct primitive characters of the same modulus $q$ in every interval of length $O_\theta(q^\theta\log T)$ above $T\gg_\theta q^\theta$, for $\theta>1/3$ (or $>1/4$ for cubefree $q$). His uniform Landau formula already uses the distance to a neighboring prime power. Srinivas \cite{srinivas} obtains discrepancies on the $\log\log T$ scale for fixed functions in the Selberg class. The dependence on conductor is central to our result. Banks and Loftus \cite{banksloftus} prove the existence of genuinely unshared zeros in sufficiently high windows for fixed characters. Eid's preprint \cite{eid} announces polylogarithmic determination for bounded order by isolating a distinguishing prime power. For Eid we used the author's indexed abstract. This leaves open a comparison of numerical exponents with the full preprint. Our estimate combines the average over nearby prime powers with a fixed separation margin for characters of all orders. The optimal height remains open.

\section{A supply of prime evaluations}
Normalize the Fourier transform by $\F g(t)=\int_{\R} g(u)e^{itu}\dd u$. We recall the positive-half-line form of the Guinand--Weil formula. For a smooth $g$ compactly supported in $(0,\infty)$,
\begin{align}
 \sum_\gamma m_\chi(\gamma)\F g(\gamma)
 &=\delta_\chi\!\int_0^\infty (e^{u/2}+e^{-u/2})g(u)\dd u
 -\int_0^\infty g(u)\frac{e^{-(a+1/2)u}}{1-e^{-2u}}\dd u \notag\\
 &\hspace{18mm}-\sum_{n\ge2}\frac{\Lambda(n)\chi(n)}{\sqrt n}\,g(\log n).
 \label{eq:weil}
\end{align}
Here $a=(1-\chi(-1))/2$ and $\delta_\chi=1$ only for the conductor-one character. In this formula and other sums indexed by $\gamma$ with an explicit multiplicity factor, each distinct ordinate occurs once. Without that factor a stated zero list counts occurrences. Applying the usual formula \cite[Lemma 9]{carneiro} to $\F g$ selects $\chi(n)g(\log n)$ on this support. The reflected term contains $g(-\log n)=0$ and the conductor term contains $g(0)=0$. The integral representation of the digamma function, with $t=2u$ in its usual exponential kernel, gives the displayed negative gamma term. The two terms for the conductor-one character are the pole contributions of the completed zeta function. Complex tests are allowed by linearity. Approximation by smooth tests on the same half-line extends \eqref{eq:weil} to our compact $C^2$ tests, with uniform integrable Fourier decay $O((1+|t|)^{-2})$ against the local count
\begin{equation}\label{eq:count}
 \sum_{v\le\gamma\le v+1}m_\chi(\gamma)\ \ll\ \log(3q(2+|v|)).
\end{equation}
For $q>1$ this follows by differencing the count for $|\gamma|\le T$ \cite{bennett}, bounding a one-sided unit interval by the corresponding annulus. The zeta case follows from its classical zero count \cite{trudgian}.

Fix a real $W\in C_c^\infty((1,2))$, with $0\le W\le1$ and $I_W=\int W(y)\dd y>0$. For every primitive character $\xi$ of conductor $f$, applying \eqref{eq:weil} to $g(u)=e^{u/2}W(e^u/X)$ gives
\begin{equation}\label{eq:smooth}
 \sum_n\Lambda(n)\xi(n) W(n/X)=\delta_\xi X I_W+O_W(\sqrt X\log(3f)).
\end{equation}
Indeed $\F g(t)=X^{1/2+it}\int W(y)y^{-1/2+it}\dd y$, with the integral $O_W((1+|t|)^{-3})$ by integration by parts. Sum \eqref{eq:count} over unit intervals. The gamma term and the smaller zeta pole term are $O_W(1)$. This proves the conductor dependence in \eqref{eq:smooth} without using an unsmoothed prime-counting error.

\begin{lemma}[Many distinguishing primes]\label{lem:primes}
For two distinct $\chi,\psi\in\mathcal C(Q)$ at least $b X/\log X$ primes $p\in[X,2X]$ satisfy $|\chi(p)-\psi(p)|\ge1$, for an absolute $b>0$. If $\psi\notin\{\chi,\bar\chi\}$, at least $b X/\log X$ primes in $[X,2X]$ satisfy $|\Re\chi(p)-\Re\psi(p)|\ge1/2$.
\end{lemma}
\begin{proof}
Write $M=\mathrm{lcm}(q_\chi,q_\psi)\le Q^2$, and view the characters on $(\Z/M\Z)^\times$. Distinct primitive characters remain distinct there by uniqueness of the primitive inducing character. Give a prime $p\nmid M$ weight $w_p=(\log p)W(p/X)$. For every character $\theta$ of this group, \eqref{eq:smooth} for its primitive representative gives
\[
 \sum_{p\nmid M}w_p\theta(p)=\boldsymbol1_{\theta=1}X I_W
 +O_W(\sqrt X\log(3M)+\sqrt X\log^2(3X)+\log(2M)).
\]
The second error bounds all higher prime powers, since there are $O(\sqrt X\log(3X))$ of them up to $2X$. Removing primes dividing $M$ costs at most $\log M$. The main term for $\theta=1$ also gives the total weight. Choosing $A$ large makes the relative error as small as desired, uniformly in $Q$. In particular $\log^2(3X)/\sqrt X$ tends uniformly to zero as $A$ increases.

On the group with uniform probability measure, the mean of $|\chi-\psi|^2$ equals 2. The mean of $(\Re\chi-\Re\psi)^2$ is at least 1 when the conjugacy pairs differ. Their real parts are orthogonal, and the mean square of each is 1 for a real character and $1/2$ otherwise. Each square is a sum of character monomials with total absolute coefficient at most 4. The prime-weighted means thus differ from the group means by at most $1/4$ after increasing $A$. The squares are bounded by 4, so a positive fraction of the total weight attains at least 1 in the first case and at least $1/4$ in the second. Now $w_p\le\log(2X)$. Decrease $b$ to cover both cases.
\end{proof}

\section{Controlling the other prime powers}
Set $h=1/8$. For $p\in[X,2X]$ and $R\ge1$ define the arithmetic quantity
\begin{equation}\label{eq:leak}
 J_p(R)=
 \sum_{\substack{n=\ell^k\ne p\\ |\log(n/p)|\le h}}
 \frac{\log\ell}{\log p}\sqrt{\frac pn}\,
 \min\!\left(1,\frac{64}{R^2|\log(n/p)|^2}\right),
\end{equation}
where $\ell$ is prime and $k\ge1$. This majorant depends only on ordinary prime powers and decreases with $R$. Its utility is that an average over all possible probe primes suffices, even though individual $J_p$ can be large.

\begin{lemma}[Averaged interference]\label{lem:sieve}
For any fixed $\varepsilon,\eta>0$, there is an absolute $B_0=B_0(\varepsilon,\eta)$ such that, for all sufficiently large $X$, all but at most $\eta X/\log X$ primes in $[X,2X]$ satisfy $J_p(B_0X/\log X)\le\varepsilon$.
\end{lemma}
\begin{proof}
For $n$ in \eqref{eq:leak} we have $X/2<n<3X$, $|\log(n/p)|\ge|n-p|/(3X)$, and the factor preceding the minimum is bounded by an absolute constant. Write $M_0=24X/R$, $V(j)=\min(1,(M_0/j)^2)$. For a positive even integer $j\le3X$, the uniform upper-bound sieve for prime pairs gives
\[
 \#\{s\le3X:s,\ s+j\text{ prime}\}\ll
 \frac{X}{\log^2 X}\,\frac j{\varphi(j)} .
\]
For example this is Ford's Corollary 2.6 \cite{ford}. It also follows by sieving out $0,-j$ for primes up to a small fixed power of $X$ \cite[Lemma 3.2]{tao}. Odd gaps do not occur between the odd primes under consideration. We have
\[
 \sum_{j\le Y}\frac j{\varphi(j)}
 =\sum_{d\le Y}\frac{\mu^2(d)}{\varphi(d)}
                 \left\lfloor\frac Yd\right\rfloor
 \le Y\prod_\ell\left(1+\frac1{\ell(\ell-1)}\right)\ll Y.
\]
Summing first below $\max(1,M_0)$ and then on doubling intervals gives
$\sum_{j\ge1} V(j)j/\varphi(j)\ll 1+M_0$. Summing the pair bound over both directions therefore bounds the contribution from $n$ prime to $\sum_pJ_p$ by $O(X(1+M_0)/\log^2 X)$.

There are $O(\sqrt X\log X)$ proper prime powers under $3X$. For each such $n$, summing $V(|p-n|)$ even over all integers $p\ne n$ costs $O(1+M_0)$. Altogether,
\begin{equation}\label{eq:sumJ}
 \sum_{X\le p\le 2X}J_p(R)
 \ll (1+X/R)\left(\frac X{\log^2 X}+\sqrt X\log X\right).
\end{equation}
At $R=B_0 X/\log X$, divide by $X/\log X$ and the upper bound is $O(1/B_0)+o(1)$ as $X\to\infty$. Choose $B_0$ and then $X$ so this is $\le\varepsilon\eta$. Markov's inequality proves the lemma.
\end{proof}

Take $\varepsilon=1/64$, $\eta=b/2$ with $b$ from Lemma~\ref{lem:primes}. Enlarge $A$ so that all values $X=A L^2$ qualify. Let $\Pp$ be the primes in $[X,2X]$ for which $J_p(B_0X/\log X)\le1/64$. It is a single set for all characters at once. Every pair in either assertion of Lemma~\ref{lem:primes} retains at least $(b/2)X/\log X$ witnesses in $\Pp$. There are $O(X/\log X)$ primes altogether, by the elementary Chebyshev upper bound.

\section{Reading a prime value from a window of zeros}
For $R$ satisfying \eqref{eq:height}, set $r=\lfloor hR/32\rfloor$ and increase the absolute constants so that $r\ge4$. Define
\[
 B_r(v)=(1-v^2/h^2)_+^r,\qquad
 K(v)=B_r(v)\sinc^2(Rv/8),\qquad \sinc t=\frac{\sin t}{t},\quad \sinc0=1 .
\]
Thus $K(0)=1$, $0\le K\le1$, and $\|K\|_1\le8\pi/R$. For each $p\in\Pp$ put
\begin{align}
 f_p(t)&=e^{it\log p}\F K(t-U),\notag\\
 Z_p(\chi;U,R)&=-\frac{\sqrt p}{\log p}
       \sum_{\gamma\in I}m_\chi(\gamma)f_p(\gamma).\label{eq:reader}
\end{align}
Everything on the right is determined by the list, the designated prime and the endpoints of the interval.

\begin{proposition}[Uniform reading error]\label{prop:read}
With the constants above and $B$ sufficiently large,
\[
 |Z_p(\chi;U,R)-\chi(p)|\le\tfrac18
 \qquad(\chi\in\mathcal C(Q),\ p\in\Pp).
\]
The same bound holds with either endpoint of $I$ omitted.
\end{proposition}
\begin{proof}
Apply \eqref{eq:weil} with $g(u)=e^{-iU(u-\log p)}K(u-\log p)$, so $\F g=f_p$. The contribution from $n=p$, after normalization as in \eqref{eq:reader}, is $\chi(p)$. Other prime powers cost at most $J_p(R)\le1/64$. On the support, $e^u\asymp X$. The absolute gamma term before normalization is $O(1/(R\sqrt X))$ and the zeta pole terms are $O(\sqrt X/R)$, uniformly in the modulation $U$. After normalization their combined contribution is $O(X/(R\log X))$, as small as desired by enlarging $B$ in \eqref{eq:height}.

It remains to control the discarded zeros. Write $B_r(v)=(1-v/h)^r(1+v/h)^r$ for $|v|<h$. The endpoint zeros through derivative order $r-1$ permit $r$ integrations by parts. In the Leibniz expansion of $B_r^{(r)}$ the product of falling factorials is at most $r^r$, and the binomial sum of the remaining powers is $((1-v/h)+(1+v/h))^r=2^r$. Thus
\[
 |\F B_r(t)|\le 2h\min\left(1,\left(\frac{2r}{h|t|}\right)^r\right).
\]
The sinc-square factor is the characteristic function of the sum of two independent uniform variables on $[-R/8,R/8]$. Hence $\F K$ is an average of translates of $\F B_r$ by at most $R/4$. For $|t|\ge R$ we obtain
\begin{equation}\label{eq:tail}
 |\F K(t)|\le2h\left(\frac R{12|t|}\right)^r.
\end{equation}
On $2^jR\le |\gamma-U|\le2^{j+1}R$, \eqref{eq:count} gives $O(2^jR(\log(3Q(3+|U|+R))+j+1))$ zeros with multiplicity. Sum \eqref{eq:tail} on those bands to obtain
\[
 \sum_{|\gamma-U|\ge R}m_\chi(\gamma)|f_p(\gamma)|
 \ll R e^{-c_0 R}\log(3Q(3+|U|+R))
\]
for some absolute $c_0>0$. To verify uniformity, put $D_0=\log(e^eQ(3+|U|))$. The last log is $O(D_0+\log(2+R))$. For sufficiently large $B$, \eqref{eq:height} gives $D_0\le e^{c_0R/4}$ and $R\gg X/\log X$. The normalized tail is then $O((\sqrt X/\log X)e^{-c_0R/2})$, less than $1/64$ after enlarging the constants. Including the endpoints in this tail estimate proves the endpoint assertion as well. Combining the errors gives the claimed margin.
\end{proof}

\paragraph{Why the height saves two logarithmic factors.}
The article reads every prime power up to $X$ with a test supported in a logarithmic interval of width $O(1/X)$. Suppressing the zero tail then leads to a height of order $X\log X$ for $X\asymp L^2$. Here the polynomial factor of $K$ has fixed support and supplies exponential decay beyond the observation window. The sinc-square factor resolves primes to an ordinary scale of order $X/R$. Once $X/R$ is a small multiple of $\log X$, the upper-bound sieve controls the resulting interference for enough primes to distinguish every pair. Thus the two jobs of the test can be estimated separately.

\section{Separation, errors and the number of probes}
\begin{proof}[Proof of Theorems \ref{th:main} and \ref{th:shift}]
For distinct characters choose $p\in\Pp$ with $|\chi(p)-\psi(p)|\ge1$. Proposition~\ref{prop:read} and the triangle inequality give
\[
 \left|\sum_{\gamma\in I}(m_\chi(\gamma)-m_\psi(\gamma))f_p(\gamma)\right|
 \ge\frac34\frac{\log p}{\sqrt p}.
\]
On the entire real line
\[
 |f_p(t)|\le\frac{8\pi}{R},\qquad |f'_p(t)|\le\frac{8\pi(\log p+h)}{R},
\]
by differentiating the test integral. Sum the first bound over the unmatched occurrences and use the second times the distance for each pair. This proves \eqref{eq:transport} and the assertion about exact matches. The proof applies to half-open intervals too.

Take $R=C_1X/\log X$ with $C_1$ large and $U=R$. This satisfies \eqref{eq:height}, and $I$ with its left endpoint omitted is $(0,2R]$. More generally the bounds apply for any $R\asymp H_Q$ with a sufficiently large lower constant and $U=R$, so take $2R=C H_Q$ with a single absolute $C$. By \eqref{eq:count}, there are $O(R\log(3Q(3+R)))=O(RL)$ occurrences per list. If every match has displacement at most $c/L^2$, then $D\ll cR/L$. Since $\sqrt X=\sqrt A\,L$, a sufficiently small absolute $c$ lets us absorb its contribution in \eqref{eq:transport}. It follows that $E\gg R\log X/\sqrt X\gg L$. This proves (i) and (ii).

For (iii), take $U=0$, $R=C H_Q$. Since $K$ is real and even, the real part of $f_p(\gamma)$ is $\cos(\gamma\log p)\F K(\gamma)$, an even function. Thus every $\Re Z_p$ can be computed from the folded list, including its multiplicities. For nonconjugate pairs some retained $p$ has $|\Re\chi(p)-\Re\psi(p)|\ge1/2$, exceeding the total approximation error $1/4$. Conversely $L(\bar s,\bar\chi)=\overline{L(s,\chi)}$, so conjugation reflects the signed list and preserves its folding.
\end{proof}

\begin{corollary}[An order-optimal number of common primes]\label{cor:dict}
There is a set of $k=O(\log(2Q))$ primes from $\Pp$, chosen independently of the unknown character, and an absolute $d>0$, such that every distinct pair has $|\chi(p)-\psi(p)|\ge1$ on at least $dk$ of them. The same set distinguishes nonconjugate pairs by their real parts, with threshold $1/2$ on at least $dk$ coordinates. A decoder succeeds with arbitrary corruptions in fewer than $dk/2$ coordinates if every other coordinate has error at most $1/8$. The order of $k$ cannot be reduced for dictionaries consisting of exact prime evaluations, even restricted to quadratic characters.
\end{corollary}
\begin{proof}
A positive fraction $\beta$ of $\Pp$ witnesses each required pair, uniformly by the two lemmas. The number of primitive characters under consideration is at most $\sum_{q\le Q}\varphi(q)\le Q^2$. Select $k=\lceil C_2\log(2Q)\rceil$ distinct primes uniformly without replacement. A hypergeometric lower-tail bound gives probability $\le\exp(-\beta k/8)$ of seeing fewer than $\beta k/2$ witnesses for any specified pair. A union bound over at most $2Q^4$ tests proves existence. If the pool has fewer than $k$ members, take the entire pool. All sets involved are finite and a successful set can also be found by enumeration. Choose a candidate agreeing within $1/8$ on the largest number of observed coordinates. On each genuine witnessing coordinate the two radius-$1/8$ disks are disjoint, also in the real-part problem. Fewer than $dk/2$ corruptions preserve unique decoding of the character for complex readings and of its conjugacy class for real readings.

For the lower bound there are $\gg Q/\log Q$ prime conductors in $(Q/2,Q]$, each furnishing a different quadratic character. A prime evaluation has at most three values for these characters. Hence $3^k\gg Q/\log Q$. For the particular primes $p\le2X<Q/2$ at large $Q$ one may replace 3 by 2. This lower bound concerns prime evaluations and leaves open the optimal height of the zero data.
\end{proof}

\section{Recovering the factors of an unlabeled product}
\begin{corollary}\label{cor:packet}
Fix $k\ge1$. If $P(s)=\prod_{j=1}^m L(s,\chi_j)$ with $0\le m\le k$ and $\chi_j\in\mathcal C(Q)$, the multiset of its positive zero ordinates up to $C_k H_Q$ determines $m$ and all its primitive factors with repetitions. Labels identifying the factor responsible for each zero are unnecessary.
\end{corollary}
\begin{proof}
Compare two different multisets of at most $k$ characters. Cancel their common factors and write $D(p)=\sum_j c_j\xi_j(p)$ with distinct primitive $\xi_j$, nonzero integer $c_j$, and $\sum|c_j|\le2k$. On the unit group of the least common multiple $M$ of their conductors, orthogonality gives $\mathbb E|D|^2=\sum c_j^2\ge1$. Each character $\xi_i\bar\xi_j$ in the expansion has primitive conductor at most $Q^2$, and $\log M\le2k\log Q$. Apply \eqref{eq:smooth} as in Lemma~\ref{lem:primes}, using $X=A_k L^2$. The sum of the absolute coefficients in the square is at most $4k^2$, and $|D|\le2k$. Thus $|D(p)|\ge1/2$ on $\gg_k X/\log X$ primes. Lemma~\ref{lem:sieve} allows us to keep one such prime with $J_p\le1/(64k)$. In Proposition~\ref{prop:read} the pole, gamma and tail bounds can be made $\le1/(64k)$ in total by increasing the constants depending on $k$. The error of the linear combination of readings is then $\le (2k)/(32k)=1/16$. Identical pooled zero lists would make this combination zero, contradicting $|D(p)|\ge1/2$. Orders add under multiplication at all points of the open critical strip, including common and multiple zeros. This also treats the empty product.
\end{proof}

\section{Finite checks}
\paragraph{A shared set of three primes.}
Take $q=101$, $X=10\,000$, $R=80\,000$ and $h=1/8$. With generator 2, each of the 100 characters of $(\Z/101\Z)^\times$ is indexed by an exponent in $\Z/100\Z$. The three primes below distinguish all 100. The last column counts the previously unresolved nonzero exponent differences distinguished at each step.
\begin{center}
\begin{tabular}{rrrr}\toprule
$p$ & $\operatorname{ind}_2 p$ & upper bound for $J_p(R)$ & differences resolved\\\midrule
11197 & 60 & 0.0420809 & 80\\
14107 & 32 & 0.0494959 & 16\\
11027 & 39 & 0.0329468 & 3\\\bottomrule
\end{tabular}
\end{center}
For a difference $a$, the chord is at least 1 precisely when $(a\operatorname{ind}_2 p\bmod 100)/100$ belongs to $[1/6,5/6]$. The integer calculation bounds \eqref{eq:leak} by enlarging the window to $[4p/5,5p/4]$, replacing the amplitude by 2 and using $|\log(n/p)|\ge|n-p|/\max(n,p)$. It enumerates every prime power in the enlarged interval and rounds each upper bound upwards. These are checks of arithmetic interference and separation at the stated parameters. They do not supply numerical values of $A,B,C$ in the general theorem.

\paragraph{A certified signed discrepancy at conductor five.}
Let $\chi$ be the quartic character modulo 5 with $\chi(2)=i$. The disk of radius $10^{-4}$ centered at $s_0=1/2+4.13290i$ contains one simple zero of $L(s,\bar\chi)$ and no zero of $L(s,\chi)$. This finite assertion holds unconditionally. The verification writes each function as the sum of its four progressions modulo 5. For each progression, sum the first 40 terms and use Euler--Maclaurin through $B_8$. For $\Re s\ge1/2$, $|s|\le5$, the remainders in the four progressions together have absolute value at most
\[
 \frac{4\cdot 5^7 (5\cdot6\cdots12)}{30\cdot 8!\cdot 7.5}\,200^{-7.5}<4\cdot10^{-12}.
\]
Interval arithmetic then gives $|L(s_0,\bar\chi)|<4.14\cdot10^{-6}$ and $|L(s_0,\chi)|>1.97$. The character partial sums have modulus at most $\sqrt2$. Partial summation bounds both functions by 26 on the disk of radius $1/4$ about $s_0$. Cauchy's estimates bound the error of a forward difference with step $10^{-4}$ and, together with its interval evaluation, give $|L'(s_0,\bar\chi)|>1.07$. On the circle of radius $10^{-4}$ the terms of degree at least two have modulus $<4.17\cdot10^{-6}$. Rouch\'e's theorem, comparing with $L'(s_0,\bar\chi)(s-s_0)$, gives one zero counted with multiplicity. For $\chi$ the total variation is $<0.011$, so the same disk contains none. The functional equation and conjugation give reflection $s\mapsto 1-\bar s$ for the zeros of each individual primitive function. The unique zero in this reflection-invariant disk is therefore on the critical line.

\paragraph{A numerical readout of the explicit formula.}
For the same two quartic characters, a normalized cubic spline supported within $0.36$ of $\log 2$ sees exactly the prime 2. The gamma terms agree. The difference of the two zero sums thus equals $-2i\log(2)/\sqrt2=-0.980258143\ldots i$. The numerical script obtains $-0.977437432i$ through height 32 and $-0.980248389i$ through height 256. It searches for sign changes on two meshes and cross-checks selected function values. Completeness of the 352 numerical roots through height 256 is unverified. The certificate in the preceding paragraph uses independent interval estimates.

\section{Further questions}
The exponent and the logarithmic denominator in the height bound may admit improvement. The proof uses a positive proportion of distinguishing primes at the classical $\log^2 Q$ scale and controls nearby Euler coefficients by their absolute values. Can cancellation among these coefficients recover enough measurements at a smaller bandwidth? Can one give a practical universal cutoff and explicit error budgets by tracking the sieve and test-function constants? For products, how rapidly must $C_k$ grow with the number of factors? Finally, can one force a zero genuinely absent from the other list at a comparable conductor-uniform height, even when multiplicities can be greater than one?

\begin{thebibliography}{99}\raggedright
\bibitem{lean} L. Sernau and collaborators, \emph{GRH for Dirichlet $L$-functions, in Lean 4}. Commit \texttt{72fead102f663d5c7a407964a13dc0cf9d977207}.
\href{https://github.com/lsernau/grh-lean/tree/72fead102f663d5c7a407964a13dc0cf9d977207}{Repository};
\href{https://github.com/lsernau/grh-lean/blob/72fead102f663d5c7a407964a13dc0cf9d977207/grh/src/Statements.lean}{statements};
\href{https://github.com/lsernau/grh-lean/blob/72fead102f663d5c7a407964a13dc0cf9d977207/grh/src/Main.lean}{exports}.
\bibitem{article} ProveIt Contributors, \emph{Beyond the Riemann hypothesis for Dirichlet $L$-functions: consequences, transfers and limits}, 26 November 2026. Commit \texttt{10802a6b360d5cdce48698261bc305e57a934b1f}.
\href{https://github.com/VladimirReshetnikov/ProveIt/tree/10802a6b360d5cdce48698261bc305e57a934b1f/Analysis/GeneralizedRiemannHypothesis/Article}{Article and provenance};
\href{https://github.com/VladimirReshetnikov/ProveIt/blob/10802a6b360d5cdce48698261bc305e57a934b1f/Analysis/GeneralizedRiemannHypothesis/Article/sections/rec.tex#L2319-L2354}{quadratic recovery};
\href{https://github.com/VladimirReshetnikov/ProveIt/blob/10802a6b360d5cdce48698261bc305e57a934b1f/Analysis/GeneralizedRiemannHypothesis/Article/sections/rec.tex#L2627-L2633}{height question}.
\bibitem{banks} W. D. Banks, Comparing zeros of distinct Dirichlet $L$-functions, \emph{Publ. Math. Debrecen} \textbf{106} (2025), 327--335. \href{https://arxiv.org/abs/2302.07073}{arXiv:2302.07073}.
\bibitem{banksloftus} W. Banks and K. Loftus, Unshared zeros of Dirichlet $L$-functions, preprint (2026). \href{https://arxiv.org/abs/2607.22930}{arXiv:2607.22930}.
\bibitem{srinivas} K. Srinivas, Distinct zeros of functions in the Selberg class, \emph{Acta Arith.} \textbf{103} (2002), 201--207. \href{https://doi.org/10.4064/aa103-3-1}{doi:10.4064/aa103-3-1}.
\bibitem{eid} S. Eid, Effective determination of primitive Dirichlet characters by their zeros up to polylogarithmic height, preprint (July 2026). \href{https://doi.org/10.5281/zenodo.21326652}{doi:10.5281/zenodo.21326652}.
\bibitem{carneiro} E. Carneiro, M. B. Milinovich, E. Quesada-Herrera and A. P. Ramos, Fourier optimization, the least quadratic non-residue, and the least prime in an arithmetic progression. \href{https://arxiv.org/abs/2404.08380}{arXiv:2404.08380}, Lemma 9. For a form including poles, see E. Carneiro, V. Chandee and M. B. Milinovich, \emph{A note on the zeros of zeta and $L$-functions}, Eq.~(4.6), \href{https://arxiv.org/abs/1503.00955}{arXiv:1503.00955}.
\bibitem{bennett} M. A. Bennett, G. Martin, K. O'Bryant and A. Rechnitzer, Counting zeros of Dirichlet $L$-functions, \emph{Math. Comp.} \textbf{90} (2021), 1455--1482. \href{https://personal.math.ubc.ca/~bennett/BeMaObRe-MathComp-2021.pdf}{Author's copy}.
\bibitem{trudgian} T. S. Trudgian, An improved upper bound for the error in the zero-counting formulae for Dirichlet $L$-functions and Dedekind zeta-functions, \emph{Math. Comp.} \textbf{84} (2015), 1439--1450. \href{https://doi.org/10.1090/S0025-5718-2014-02898-6}{doi:10.1090/S0025-5718-2014-02898-6}.
\bibitem{ford} K. Ford, \emph{Sieve Methods}, lecture notes, Corollary 2.6. \href{https://ford126.web.illinois.edu/sieve2023.pdf}{Author's notes}.
\bibitem{tao} T. Tao and J. Ter\"av\"ainen, The Hardy--Littlewood--Chowla conjecture in the presence of a Siegel zero, \emph{J. London Math. Soc.} \textbf{106} (2022), 3317--3378. \href{https://doi.org/10.1112/jlms.12663}{doi:10.1112/jlms.12663}, Lemma 3.2.
\end{thebibliography}
\end{document}
